\chapter{\abstractname}

Recovering camera poses and intrinsics from casual, uncalibrated video is a prerequisite for 3D reconstruction, novel view synthesis, and autonomous navigation. Strong pretrained models exist for individual sub-tasks---self-supervised pose estimation, monocular depth, optical flow, single-image calibration---but they operate in isolation, with no mechanism for joint reasoning or temporal consistency.

This thesis proposes a framework for fusing a pretrained calibration network with a self-supervised pose estimator, introducing multi-frame calibration consistency through a learned feature aggregation module. The core contribution is the Multi-Frame Calibration Transformer (MCT), a lightweight module that aggregates spatial features from a calibration backbone across video frames via cross-frame self-attention before decoding, preserving geometric structure that scalar averaging would discard. The MCT is architecture-agnostic: it inserts between any compatible feature extractor and decoder, adding only 25M trainable parameters while all pretrained components remain frozen (7.5\% of total parameters trained). Multi-frame consistency is enforced efficiently through flow composition, reducing the required computations from quadratic to linear in the number of frames.

The framework is demonstrated by fusing AnyCam (a self-supervised pose estimator) with AnyCalib (a single-image calibration network). Evaluation on MPI Sintel, TUM-RGBD, and KITTI shows that the fused system improves upon both individual models: up to 39.5\% translation direction error reduction over AnyCam and up to 32.5\% calibration MAPE reduction over per-frame AnyCalib averaging, while maintaining fully self-supervised training with no ground-truth labels.
